\section{The square of the CZX \texorpdfstring{$\Z_2$}{Z2} symmetry}%
\label{sec:asymex_czx}

On a periodic chain of $N>0$ qubits, the undecorated CZX operator acts on
computational basis vectors by
\begin{align}
	U_N\ket{t} & =(-1)^{q(t)}\ket{\bar t}.
	\label{eq:asymex_czx_basis}
\end{align}
Here $t=(t_0,\ldots,t_{N-1})\in\{0,1\}^N$, $\bar t_j=1-t_j$, and
all site subscripts are taken modulo $N$. The cyclic phase exponent is
\begin{align}
	q(t) & =\sum_{j=0}^{N-1}t_jt_{j+1}.
	\label{eq:asymex_czx_exponent}
\end{align}
Define the diagonal operator $D_N$ by
$D_N\ket{t}=(-1)^{q(t)}\ket{t}$, and let $X\ket{a}=\ket{1-a}$.
With operator products acting from right to left,
\begin{align}
	U_N & =X^{\otimes N}D_N.
	\label{eq:asymex_czx_operator}
\end{align}
For $N\geq3$, $D_N$ is the product of controlled-$Z$ gates on the
cyclic nearest-neighbor bonds. Formula~(\ref{eq:asymex_czx_exponent})
also specifies the short rings: it includes a self-loop for $N=1$ and
counts the two-site bond twice for $N=2$.

The CZX example in~\cite{Cirac2021Matrix} presents a bond-two virtual
MPO symmetry in the discussion of two-dimensional symmetry-protected
topological states. Its displayed gate order is $D_NX^{\otimes N}$.
The tensor considered here instead gives~(\ref{eq:asymex_czx_operator});
the two orders differ by $(-1)^N$. The operator with additional Pauli
$Z$ factors in~\cite[Section~III.D]{GarreRubioSchuch2024DomainWall}
is a different representative. We first compute the periodic operator
of the specified tensor, then examine the compression of its stacked
square. These calculations do not determine a three-cocycle class.

\begin{definition}[The bond-two symmetry tensor and its square]
	\label{def:asymex_czx_tensors}
	\lean{CZXCompression.czxTensor, CZXCompression.czxMPS,
		CZXCompression.identityTensor, CZXCompression.identityMPS,
		CZXCompression.czxSquare, CZXCompression.czxSquareTarget}
	\leanok
	\uses{def:mpo_tensor, def:mpo_to_mps, def:mpdo_operator_product}
	The symmetry tensor is $M^{ij}=\delta_{i,1\oplus j}T_j$, where
	$i,j\in\{0,1\}$ are respectively the physical output and input,
	$\delta_{a,b}$ is the Kronecker delta, $\oplus$ denotes addition
	modulo two, and
	\begin{align}
		T_0 & =\begin{pmatrix}1&0\\1&0\end{pmatrix}.
		\notag
	\end{align}
	\begin{align}
		T_1 & =\begin{pmatrix}0&1\\0&-1\end{pmatrix}.
		\notag
	\end{align}
	The identity tensor is the bond-one tensor
	$\delta^{ij}=\delta_{i,j}$. Read over the alphabet of index pairs,
	encoded by $2i+j$, the symmetry tensor has the four matrices
	$0$, $T_1$, $T_0$, $0$ in the order $(0,0)$, $(0,1)$, $(1,0)$, $(1,1)$.
	The stacked product has bond dimension four and matrices
	\begin{align}
		B^{ij} & =\sum_{m=0}^{1}M^{im}\otimes M^{mj}.
		\label{eq:asymex_czx_stack}
	\end{align}
	The single target for its compression is $-\delta$.
\end{definition}

The periodic boundary is the ordinary, unnormalized trace. Thus the
operator generated by $M$ has matrix coefficients
\begin{align}
	U_N(s,t) & =\tr(M^{s_0t_0}\cdots M^{s_{N-1}t_{N-1}}).
	\label{eq:asymex_czx_periodic}
\end{align}
There are no open boundary vectors. For bond indices
$\alpha,\beta\in\{0,1\}$, the local entries are
\begin{align}
	M^{ij}_{\alpha,\beta}
	 & =\delta_{i,1-j}\delta_{\beta,j}(-1)^{\alpha j}.
	\label{eq:asymex_czx_entries}
\end{align}

\begin{theorem}[Periodic computational basis kernel]
	\label{thm:asymex_czx_kernel}
	\lean{CZXCompression.mpo_czxTensor_apply, CZXCompression.mpo_czxTensor}
	\leanok
	\uses{def:asymex_czx_tensors}
	For every $N>0$ and every $s,t\in\{0,1\}^N$, the periodic
	contraction satisfies
	\begin{align}
		U_N(s,t) & =\delta_{s,\bar t}(-1)^{q(t)}.
		\label{eq:asymex_czx_kernel}
	\end{align}
	In particular, its basis action is~(\ref{eq:asymex_czx_basis}).
\end{theorem}

\begin{proof}
	\leanok
	Expand the trace over cyclic bond configurations $g$, using
	(\ref{eq:asymex_czx_entries}):
	\begin{align}
		U_N(s,t)
		 & =\sum_{g\in\{0,1\}^N}
		\prod_{j=0}^{N-1}
		\delta_{s_j,1-t_j}\delta_{g_{j+1},t_j}
		                  (-1)^{g_jt_j}.
		\label{eq:asymex_czx_contraction}
	\end{align}
	The bond constraints force $g_j=t_{j-1}$, so exactly one bond
	configuration can contribute. Its contribution vanishes unless
	$s=\bar t$. Otherwise it equals
	$(-1)^{\sum_jt_{j-1}t_j}=(-1)^{q(t)}$, by cyclic reindexing.
\end{proof}

\begin{theorem}[Positive-length unitarity]
	\label{thm:asymex_czx_unitarity}
	\lean{CZXCompression.mpo_czxTensor_mem_unitaryGroup,
		CZXCompression.czxTensor_isMPUPos}
	\leanok
	\uses{def:asymex_czx_tensors}
	The operator $U_N$ is unitary for every $N>0$.
\end{theorem}

\begin{proof}
	\leanok
	\uses{thm:asymex_czx_kernel}
	By Theorem~\ref{thm:asymex_czx_kernel}, $U_N$ is a monomial matrix.
	The map $t\mapsto\bar t$ is a permutation and every nonzero entry
	has modulus one, since
	$\overline{(-1)^{q(t)}}(-1)^{q(t)}=1$. Its columns and rows are
	therefore orthonormal.
\end{proof}

\begin{theorem}[Square of the periodic operator]
	\label{thm:asymex_czx_square_operator}
	\lean{CZXCompression.mpo_czxTensor_mul_self,
		CZXCompression.mpo_mulTensor_czxTensor}
	\leanok
	\uses{def:asymex_czx_tensors}
	For every $N>0$,
	\begin{align}
		U_N^2 & =(-1)^N\Id.
		\label{eq:asymex_czx_square_operator}
	\end{align}
	The periodic operator generated by the stacked tensor $B$ is
	also $(-1)^N\Id$.
\end{theorem}

\begin{proof}
	\leanok
	\uses{thm:asymex_czx_kernel}
	For bits $a,b$, direct evaluation gives
	$(-1)^{(1-a)(1-b)}(-1)^{ab}
		=-(-1)^a(-1)^b$. Multiplying this identity around the ring,
	each single-site sign occurs twice, so
	\begin{align}
		(-1)^{q(\bar t)}(-1)^{q(t)}
		 & =(-1)^N\prod_{j=0}^{N-1}(-1)^{2t_j}=(-1)^N.
		\label{eq:asymex_czx_phase_square}
	\end{align}
	Apply~(\ref{eq:asymex_czx_kernel}) twice and use
	$\overline{\bar t}=t$ to obtain~(\ref{eq:asymex_czx_square_operator}).
	The periodic contraction of the stacked product is the product
	of the two periodic operators, by expansion over the intermediate
	physical indices in~(\ref{eq:asymex_czx_stack}).
\end{proof}

Unitarity and~(\ref{eq:asymex_czx_square_operator}) give the inverse
and physical adjoint:
\begin{align}
	U_N^{-1} & =U_N^\dagger=(-1)^NU_N.
	\label{eq:asymex_czx_adjoint}
\end{align}
Consequently $U_N^4=\Id$. On even positive rings, $U_N$ is a
self-adjoint involution; on odd rings, $U_N^2=-\Id$.
The phase identity~(\ref{eq:asymex_czx_phase_square}) also gives
$D_NX^{\otimes N}=(-1)^NU_N$, explaining the gate-order difference.

At length one, $q(t)=t_0^2=t_0$, hence $D_1=Z$ and $U_1=XZ$,
where $Z\ket{a}=(-1)^a\ket{a}$. At length two,
$q(t)=2t_0t_1$, hence $D_2=\Id$ and $U_2=X\otimes X$.
These contractions must not be replaced by a convention that omits the
self-loop or counts the two-site edge only once.

At length zero the trace convention gives the scalar
$\tr(\Id_2)=2$, not a unitary operator. The stacked tensor similarly
gives $\tr(\Id_4)=4$, whereas the bond-one target gives
$\tr(\Id_1)=1$. All periodic operator and trace comparisons below
therefore retain their positive-length or nonempty-word hypotheses.

The physical square law supplies the periodic comparison for
asymmetric compression: $B$ has the nonempty-word traces of the
bond-one identity tensor with weight $-1$. This comparison lies
beyond the unit-weight setting. The sitewise theory
of~\cite[Appendix~A]{GarreRubio2023MPOSym} imposes a closedness
requirement for arbitrary boundary conditions. The following example
distinguishes a word-level compression from a sitewise intertwiner.

\begin{center}
	% Formula: B^{ij}=\sum_mM^{im}\otimes M^{mj}, the stacked product of
	% Definition~\ref{def:asymex_czx_tensors}.
	% Source: Definition~\ref{def:asymex_czx_tensors}.
	% Ink-to-index map: both rows are the bond-two symmetry tensor $M$;
	% the shared vertical pairleg at each site is the summed index $m$.
	% Contraction set: the vertical pairlegs (index $m$) at every site
	% and the horizontal virtual bonds within each row; the physical
	% legs $(i,j)$ remain open.
	% Boundary signature: (virtual-west=0 | virtual-east=0 |
	% physical-up=3 | physical-down=3) on both sides.
	\tenkzkernel
	\begin{tenkz}[rows={op,op}, cols=4, boundary=periodic]
		\tn[skin=mpo, ports={90:physical:$i_1$, 270:physical}]{M} &
		\tn[skin=mpo, ports={90:physical:$i_2$, 270:physical}]{M} &
		\tn[skin=dots, ports={180:virtual, 0:virtual}]{} &
		\tn[skin=mpo, ports={90:physical:$i_N$, 270:physical}]{M} \\
		\tn[skin=mpo, ports={90:physical, 270:physical:$j_1$}]{M} &
		\tn[skin=mpo, ports={90:physical, 270:physical:$j_2$}]{M} &
		\tn[skin=dots, ports={180:virtual, 0:virtual}]{} &
		\tn[skin=mpo, ports={90:physical, 270:physical:$j_N$}]{M}
	\end{tenkz}
	\quad $=$ \quad
	\begin{tenkz}[rows={op}, cols=4, boundary=periodic]
		\tn[skin=mpo, ports={90:physical:$i_1$, 270:physical:$j_1$}]{B} &
		\tn[skin=mpo, ports={90:physical:$i_2$, 270:physical:$j_2$}]{B} &
		\tn[skin=dots, ports={180:virtual, 0:virtual}]{} &
		\tn[skin=mpo, ports={90:physical:$i_N$, 270:physical:$j_N$}]{B}
	\end{tenkz}
\end{center}

For a word $w=((i_1,j_1),\ldots,(i_{|w|},j_{|w|}))$,
write $B^w=B^{i_1j_1}\cdots B^{i_{|w|}j_{|w|}}$, and use the
same convention for the other tensors. The empty word evaluates to
the identity matrix on the corresponding bond space.

\begin{theorem}[Normality of the symmetry tensor and of the target]
	\label{thm:asymex_czx_normal}
	\lean{CZXCompression.czxMPS_isNormal, CZXCompression.czxSquareTarget_isNormal}
	\leanok
	\uses{def:asymex_czx_tensors, def:normal}
	The bond-two symmetry tensor is normal: its length-two words span
	$\MN{2}$. The target $-\delta$ is normal: its bond dimension is one and
	one of its matrices is the nonzero scalar $-1$.
\end{theorem}

\begin{proof}
	\leanok
	\uses{thm:asymex_normal_units}
	Let $E_{ab}$ denote the matrix units of $\MN{2}$. Direct multiplication
	gives
	\begin{align}
		E_{00} & =\tfrac12(T_0T_0+T_1T_0).
		\notag
	\end{align}
	\begin{align}
		E_{10} & =\tfrac12(T_0T_0-T_1T_0).
		\notag
	\end{align}
	\begin{align}
		E_{01} & =\tfrac12(T_0T_1-T_1T_1).
		\notag
	\end{align}
	\begin{align}
		E_{11} & =\tfrac12(T_0T_1+T_1T_1).
		\notag
	\end{align}
	Thus every matrix unit lies in the length-two word span. For the
	target, its matrix at the letter $(0,0)$ is the nonzero scalar $-1$,
	so Theorem~\ref{thm:asymex_normal_units} applies with the single
	position of the one-by-one matrix algebra.
\end{proof}

This normality is a statement about words spanning virtual matrices;
it does not impose a normalization of the transfer map's spectral radius.

\begin{theorem}[Compression of the squared symmetry]
	\label{thm:asymex_czx_compression}
	% Principal declaration: CZXCompression.czxSquare_isReduction.
	\lean{CZXCompression.czxSquare_compression, CZXCompression.czxSquareLeft,
		CZXCompression.czxSquareRight, CZXCompression.czxSquare_isReduction,
		CZXCompression.czxSquareLeft_mul_right, CZXCompression.czxSquare_dim_eq,
		CZXCompression.czxSquare_evalWord_remainder_eq_zero}
	\leanok
	\uses{def:asymex_czx_tensors, thm:asym_multiblock}
	The stacked product of Definition~\ref{def:asymex_czx_tensors} carries
	the block triangular form of Theorem~\ref{thm:asym_multiblock} for the
	single slot $-\delta$ with $z=3$ zero slots and $4=1+3$. The
	compression pair is
	\begin{align}
		W & =(0,0,-1,0).
		\notag
	\end{align}
	\begin{align}
		V & =(1,1,-1,-1)^{\mathsf T}.
		\notag
	\end{align}
	It satisfies $WV=1$ and $WB^wV=(-\delta)^w$ for every word.
	The remainder $R^{ij}=B^{ij}-V(-\delta)^{ij}W$ satisfies
	\begin{align}
		R^{i_1j_1}R^{i_2j_2}R^{i_3j_3}R^{i_4j_4} & =0
		\notag
	\end{align}
	for every choice of four index pairs.
\end{theorem}

\begin{proof}
	\leanok
	\uses{thm:asymex_explicit_gauge}
	Set $k=(1,0,0,-1)^{\mathsf T}$ and
	$v=(1,1,-1,-1)^{\mathsf T}$. The only nonzero letters of $B$ are
	\begin{align}
		B^{00} & =\begin{pmatrix}
			          0 & 0 & 1  & 0 \\
			          0 & 0 & 1  & 0 \\
			          0 & 0 & -1 & 0 \\
			          0 & 0 & -1 & 0
		          \end{pmatrix}.
		\label{eq:asymex_czx_bzero}
	\end{align}
	\begin{align}
		B^{11} & =\begin{pmatrix}
			          0 & 1  & 0 & 0 \\
			          0 & -1 & 0 & 0 \\
			          0 & 1  & 0 & 0 \\
			          0 & -1 & 0 & 0
		          \end{pmatrix}.
		\label{eq:asymex_czx_bone}
	\end{align}
	Both annihilate $k$, their images lie in $\spn_{\C}\{k,v\}$,
	and $B^{00}v=-v$, $B^{11}v=2k-v$. Choose a basis beginning
	with $k,v$ and complete it with two further vectors. Every letter
	is upper triangular in this basis, with diagonal
	$(0,-\delta_{i,j},0,0)$. A completion compatible with the displayed
	row $W$ gives the compression pair by
	Theorem~\ref{thm:asymex_explicit_gauge} applied to the second
	coordinate. The vanishing of a product of four remainder matrices is
	clause (vi) of Theorem~\ref{thm:asym_multiblock} with $|S|+z=4$.
\end{proof}

\begin{theorem}[Periodic data of the squared symmetry]
	\label{thm:asymex_czx_trace}
	\lean{CZXCompression.czxSquare_trace_evalWord}
	\leanok
	\uses{def:asymex_czx_tensors}
	For every nonempty word $w$ over the alphabet of index pairs,
	$\tr(B^w)=(-1)^{|w|}\tr(\delta^w)$. This is
	$U_N^2=(-1)^N\Id$ at the level of periodic coefficients.
\end{theorem}

\begin{proof}
	\leanok
	\uses{thm:asymex_czx_compression, prop:asym_compression_trace}
	Proposition~\ref{prop:asym_compression_trace} applied to the compression
	of Theorem~\ref{thm:asymex_czx_compression} gives
	$\tr(B^w)=\tr((-\delta)^w)$. Pulling the scalar out of
	each letter contributes $(-1)^{|w|}$.
\end{proof}

\begin{theorem}[Both sitewise intertwiner spaces vanish]
	\label{thm:asymex_czx_no_intertwiner}
	\lean{CZXCompression.czxSquare_right_intertwiner_eq_zero,
		CZXCompression.czxSquare_left_intertwiner_eq_zero}
	\leanok
	\uses{def:asymex_czx_tensors}
	If $v\in\C^4$ satisfies $B^{ij}v=(-\delta)^{ij}v$ for every index pair,
	then $v=0$. If a row vector $u\in\C^4$ satisfies
	$uB^{ij}=(-\delta)^{ij}u$ for every index pair, then $u=0$.
\end{theorem}

\begin{proof}
	\leanok
	Use the two nonzero matrices~(\ref{eq:asymex_czx_bzero})
	and~(\ref{eq:asymex_czx_bone}). The right-intertwiner equations are
	\begin{align}
		(v_2,v_2,-v_2,-v_2)^{\mathsf T}
		 & =-(v_0,v_1,v_2,v_3)^{\mathsf T},
		\notag
	\end{align}
	and
	\begin{align}
		(v_1,-v_1,v_1,-v_1)^{\mathsf T}
		 & =-(v_0,v_1,v_2,v_3)^{\mathsf T}.
		\notag
	\end{align}
	They give $v_1=v_2=0$ and then $v_0=v_3=0$.
	For a left intertwiner, the equations are
	\begin{align}
		(0,0,u_0+u_1-u_2-u_3,0) & =-u,
		\notag
	\end{align}
	and
	\begin{align}
		(0,u_0-u_1+u_2-u_3,0,0) & =-u.
		\notag
	\end{align}
	The equation for $B^{00}$ gives $u_0=u_1=u_3=0$,
	and the equation for $B^{11}$ gives $u_2=0$.
\end{proof}

\begin{remark}[Failure of closure under conjugate transposition]
	\label{rem:asymex_czx_anomaly}
	\uses{cor:asym_anomaly_not_star}
	Corollary~\ref{cor:asym_anomaly_not_star} applies to
	Theorem~\ref{thm:asymex_czx_no_intertwiner}: the algebra generated by
	the matrices of $B$ is not closed under conjugate transposition in
	these coordinates. This coexists with physical unitarity, proved in
	Theorem~\ref{thm:asymex_czx_unitarity}. The absence of a sitewise
	intertwiner persists under an invertible change of bond coordinates;
	closure under conjugate transposition refers to the specified inner
	product. Neither this failure of closure nor nonsplitting identifies
	an anomaly class. Such a claim requires coherent local fusion data
	and a separate calculation of the fusion-tensor associator and its
	cohomology class.
\end{remark}
